\begin{abstract}
Small software projects rarely close the loop with the people who report bugs or request features. Capture is easy; what follows---acknowledgement, a first substantive response, a decision, and notice of the outcome---is unowned, unmeasured, and therefore slow. This paper presents the Fast Feedback Resolution System (\ffrs), a model that runs each feedback item through a five-stage pipeline (capture, acknowledge, route, respond, close) in which every stage is a timestamp recorded by the system that owns it. Its defining element is an \emph{agentic Respond stage}: a scheduled AI agent produces the first substantive response to every item---for changes it can make, a complete pull request with tests and evidence for a developer to review and merge; otherwise a proposal that is executed only after the requester accepts it and a reviewer confirms it. No agent action takes effect without a human decision. Because the stages are timestamps, the metrics follow without manual bookkeeping. We describe a dependency-minimal, tool-agnostic reference architecture---a 5\,KB (gzipped) widget, one serverless function, the project's existing issue tracker as system of record, a private sidecar for the only personal datum, and a scheduled coding agent---and its open implementation, run as one shared, multi-tenant service and evaluated on two sites from 23~September~2026: the project's reference site, whose Respond stage is agentic, and an anonymous B2B quick-commerce platform that responds manually through the same pipeline. We report \todo{$n$} items over \todo{$k$} weeks (median time to first response \todo{$x$}, loop closure \todo{$y$}\%, agent share \todo{$z$}\%) against pre-registered targets and characterise when a minimal, human-reviewed agentic loop is sufficient. All artefacts are open; the reference site's metrics are recomputable from public data.
\end{abstract}

\begin{IEEEkeywords}
user feedback, feedback loops, AI agents, human-in-the-loop, issue tracking, DevOps metrics, software maintenance
\end{IEEEkeywords}
